
\titre{ {} }
\theme{Division euclidienne et polynôme}
%\isuuid{ba2u}
\auteur{Quentin LIARD}
%\datecreate{2025-02-13}
\organisation{AMSCC}

\contenu{On pose le polynôme $P$ défini par $$P(X)=2X^5-4X^4-2X^3+3X^2-5X-4.$$
 

 \begin{enumerate}
	\item \question{Déterminer le quotient $Q$ et le reste $R$ de la division euclidienne de $P$ par $X^2-2X-1$.}
	%\indication{}
    \reponse{$P(X)=(X^2-2X-1)Q(X)+X-1.$}
    \item \question{Déterminer les racines du polynôme $X^2-2X-1$.}
   % 	\indication{}
    \reponse{Le discriminant est négatif et égal à $8$, on a ainsi $z_1=1-\sqrt{2}$ et $z_2=1+\sqrt{2}$}
    \item \question{En déduire $P(\sqrt{2}+1).$}
    \reponse{$P(\sqrt{2})=\sqrt{2})$.}
\end{enumerate}

}